\documentclass[10pt,a4paper]{amsart}
\usepackage[margin=19mm]{geometry}
\usepackage{amsmath,amssymb,mathtools}
\usepackage[scr=rsfs,cal=euler]{mathalfa}
\usepackage{xfrac}
\usepackage[unicode,hidelinks,bookmarks=false]{hyperref}
\newcommand{\ie}{i.e.\ }
\newcommand{\cf}{cf.\ }
\newcommand{\resp}{respectively\ }
\newcommand{\Subsection}{\subsection}
\providecommand{\nobreakdash}{\nobreak\hbox{-}}
\setlength{\emergencystretch}{2em}
\multlinegap0pt
\usepackage{amssymb}
\usepackage{xcolor}
\usepackage{algorithm}\usepackage{algpseudocode}
\newtheorem{proposition}{Proposition}
\newtheorem{theorem}{Theorem}
\title{An Algebraic Approach to the Fundamental Theorem of Algebra}
\author{Priyabrata Mandal, Sajad A. Sheikh}
\date{Source: arXiv:2608.19256v1 (CC BY 4.0)}
\begin{document}
\maketitle

\noindent\emph{Source and licence.} An Algebraic Approach to the Fundamental Theorem of Algebra. Version arXiv:2608.19256v1 (CC BY 4.0), distributed by arXiv under the Creative Commons Attribution 4.0 International licence, http://creativecommons.org/licenses/by/4.0/, as recorded in the arXiv licence metadata for this exact version and shown on the abstract page https://arxiv.org/abs/2608.19256v1. Full licence notice: \texttt{licenses/arxiv-2608.19256-CC-BY-4.0.txt}.\par\smallskip

\noindent\textit{Editorial scope.} Counted unit: the article's theorem closed (source0.tex lines 232--234). The article's complete original proof of it is delivered with it (source0.tex lines 235--251, one \texttt{proof} environment of 17 lines). No mathematics is written, repaired or re-derived: the statement and the proof are the article's own lines, in the article's own notation, and no retained line is substituted or re-flowed.  Retained uncounted as the source-local context the proof needs: lines 184--186.  Nothing was omitted from any retained span, so the package carries no omission record.  The retained lines cite 1 works, whose entries are transcribed verbatim from the article's own bibliography source in the e-print and delivered with short editorial labels recorded in the item's reference mapping.  No retained line contains a figure environment, an image-inclusion command or drawing code, so no image is delivered and the package declares no asset dependency.  Editorial formatting only, in addition: an AMS \texttt{amsart} preamble that restates the article's own declarations and macros used by the retained lines, namely the environments \texttt{proposition}, \texttt{theorem} on separate counters, together with the standard packages those lines need.  The retained environments print in this document as Proposition 1, Theorem 1 in order of appearance, whereas the article prints the same items with the numbering of its own full document.  The complete native source of the article as distributed in the arXiv e-print for this exact version is bundled unchanged as \texttt{sources/208074/source0.tex}.
\begin{proposition}[\cite{lam}, Chapter 8, Proposition 1.6]\label{Euclidean is Pythagorean}
    A field $\mathbb E$ is Euclidean if and only if $\mathbb E$ is a Pythagorean field with a unique ordering.
\end{proposition}

\begin{theorem}\label{Euclidean is quad closed}
    Let $\mathbb E$ be a Euclidean field. Then $\mathbb E(\sqrt{-1})$ is quadratically closed. 
\end{theorem}

\begin{proof}
    Consider an arbitrary non-zero element $a+(\sqrt{-1})b \in \mathbb E(\sqrt{-1})$. To prove that $\mathbb E(\sqrt{-1})$ is quadratically closed, it suffices to show that there exists $x,y \in \mathbb E$ such that  $a+(\sqrt{-1})b=(x+(\sqrt{-1})y)^2$. Expanding the square and comparing both sides, we get $a=x^2-y^2$ and $b=2xy$. Therefore $y=\frac{b}{2x}$. Substituting the value of $y$, we get
\begin{align*}
    & a=x^2- \frac{b^2}{4x^2}\\ 
     \text{i.e., } &x^4-ax^2-\frac{b^2}{4}=0
\end{align*}
    Solving, we get $x^2=\frac{1}{2}(a+\sqrt{a^2+b^2})$. As $\mathbb E$ is Pythagorean (by Proposition \ref{Euclidean is Pythagorean}), we have $a^2+b^2=c^2$ for some $c \in \mathbb E$. Thus, $x^2=\frac{1}{2}(a \pm c)$. 
    \medskip 
    
    \textbf{Case I:} If $a\geq c$, then $a \pm c \geq 0$. Since $E$ is Euclidean, we have $a \pm c =d^2$ for some $d \in \mathbb E$. Thus $x^2=\frac{1}{2}d^2$. Hence, $x= \pm \frac{1}{\sqrt{2}}d \in \mathbb E$ and $y=\frac{b}{2x} \in \mathbb E$.
\medskip 

    \textbf{Case II:} Suppose $a \pm c \leq 0$. Then we have $a \pm c = -d^2$ for some $d \in \mathbb E$. Thus  $x^2=-\frac{1}{2}d^2$. Hence, $x= \pm \frac{1}{\sqrt{2}}(\sqrt{-1})d \in \mathbb E(\sqrt{-1})$ and $y=\frac{b}{2x} \in \mathbb E(\sqrt{-1})$.
    
   % and so $x=\frac{1}{\sqrt{2}}(a+\sqrt{a^2+b^2})^\frac{1}{2}$. Hence, $y=\frac{b}{\sqrt{2}}(a+\sqrt{a^2+b^2})^{-\frac{1}{2}}$.
    Therefore, every element of $\mathbb E(\sqrt{-1})$ is a square. Thus, $\mathbb E(\sqrt{-1})$ is quadratically closed.
\end{proof}

\begin{thebibliography}{99}
\bibitem[lam]{lam} T. Y. Lam: \textit{Introduction to Quadratic Forms over Fields}. Graduate Studies in Mathematics, Volume 67.
\end{thebibliography}
\end{document}
